This draft file contains some possible directions to go in. Non of them are finished yet. 



\begin{itemize}
  \item For $S = \lim S_n$, we can consider the free Boolean algebra on $S_n$ and quotient out the relations making the induced map from $S$ into a morphism of Boolean algebras. 
    It is not generally true that those quotients are finite though, and one has to maybe do some clever stuff from there. 
  \item (After conversation with Evan): If the above can be made to work, we can study a more general version of the theorem.  One can consider the free algebra-forgetful functor adjunction. $\begin{tikzcd}
      BA 
      \arrow[r, bend right = -20 ,"F"{name=2}]
      &
      Set 
      \arrow[l,"U"{name = 1},bend right = -20]
      \arrow[to=2, from=1,draw=none,"\dashv"{anchor=center,rotate=-90}]
    \end{tikzcd}$. Then for any Boolean algebra $B$ we have an adjunction 
$\begin{tikzcd}
      {UB/Set}
      \arrow[r,""{anchor=center,name = a},bend right = -20]
      &
      {B/BA}
      \arrow[l, bend right = -20 ,"B/U"{name=b,anchor=center}]
      \arrow[to=a, from=b,draw=none,"\dashv"{anchor=center,rotate=-90}]
    \end{tikzcd}$
    where the left adjoint is given by the above construction. One could find out in general whether the left adjoint preserves sequential limits. 
  \item (After asking Chaitanya about the above), the general statement would have something to do with literature on lifting pro-objects along monadic functors. 
  \item (After conversation with Lorenzo) One thing we haven't looked at yet is what it means for the underlying boolean algebra of $S$. $2^S$ now has both a boolean algebra structure and a co-boolean algebra structure. Note that it has this as a set, not as an object of the category of Boolean algebras (it's boolean algebra structure does not contain of Boolean morphisms). 
  \item (After conversation with Jonas) It is known that if $B$ is a Stone ring, then we can consider the poset of open ideals $I\subseteq B$ and consider $B/I$ which are all finite sets, and we can show that the limit of $B/I$ is $B$. 
\end{itemize}

\begin{remark}
  It is not the case that a Stone space with decidable equality is finite. For example, any closed proposition is Stone and has decidable equality, but not any closed proposition is finite. 

  It is the case that if $S$ is Stone and has decidable equality, $S$ embeds into a finite set. Consider that if $S$ has decidable equality we have the singleton map $S\hookrightarrow 2^S$, which by Scott factors over some finite set. 
\end{remark}

\begin{theorem}
  Consider $\mathbb T$ any finite algebraic theory. Let $S$ be a Stone type with a $\mathbb T$-structure. Then $S$ can be written as a limit of finite types that have a $\mathbb T$-structure also. 
\end{theorem}
\begin{proof}
  We will first write $S$ as a limit of 
\end{proof}

The following came out of a chatgpt conversation
\begin{lemma}
  Let $\sim$ be a closed relation on $2$. 
  Then $2/\sim$ is totally disconnected.
\end{lemma}
We have two proofs for this:
\begin{proof}
  Denote $\pi : 2 \to 2/\sim$ for the quotient map. 
  We need to show that for any $x,y:2/\sim$ we have that $\Pi_{D:2/\sim \to 2} D(x) = D(y) \to x = y$. It suffices to show this for $x= \pi(0), y = \pi(1)$.
  So assume that $D(\pi(0)) = D(\pi(1))$ for all $D: 2/\sim \to 2$, we shall show that $0\sim 1$. 
As $\sim$ was closed, we can show the double negation. 
  Suppose that $0\nsim 1$, then $E: 2\to 2$ given by $a \mapsto a = 0$ respects $\sim$, and thus induces a map $2/\sim\to 2$, and thus $E(0) = E(1)$, contradicting that $0 \neq 1$. 
\end{proof}
\begin{proof}
  Denote $P = 0 \sim 1$. 
  We claim that $2/\sim = \neg P \to 2$. 
  The latter is Stone by Tychonov, so this would be sufficient. 
  For the map $2/\sim \to (\neg P \to 2)$, consider that if $\neg P$, then $2/\sim = 2$. 

  For the map $\neg P \to 2\to 2/\sim$, let $f:\neg P \to 2$. 
  We claim that 
  $(\Pi_{x : \neg P} fx = 0) \vee (\Pi_{x : \neg P} fx = 1)$. By Tychonov, this is closed, hence double negation stable. 
  Then sending the left case to $0$ and the right case to $1$ defines a map
  $$
  (\Pi_{x : \neg P} fx = 0) + (\Pi_{x : \neg P} fx = 1) \to 2
  $$
  that induces the desired map
  $$
  (\Pi_{x : \neg P} fx = 0) \vee (\Pi_{x : \neg P} fx = 1) \to 2/\sim.
  $$
\end{proof}

The first proof generalizes as follows:
\begin{lemma}
  Let $F$ finite and $\sim $ a double-negation stable equivalence relation on $F$. 
  Then $X:=F/\sim$ is totally disconnected. 
\end{lemma}
\begin{proof}
  We need to show 
  $$
  \Pi_{x,y:X} ((\Pi_{D:X\to 2} D(x) = D(y)) \to x = y)
  $$
  Denote $\pi: F\to X$ for the quotient map. Consider $a,b:F$ and suppose $\Pi_{D:X\to 2} D(\pi(a)) = D(\pi (b))$. We will show that $\neg \neg (\pi(a) = \pi(b))$. 
  Suppose $\pi(a)\neg \pi(b)$, we will show a contradiction. As we have $\neg \neg (a \sim c \vee a\nsim c$ for all $c:F$ and $F$ is finite, we have $\neg\neg \Pi_{c:F} (a\sim c \vee a \nsim c)$. As we're showing a contradiction, we may thus assume $\sim$ is decidable. But then we can define $E: F\to 2$ by $E(x)=a \sim x$, and this induces a map $D:X \to 2$ seperating $\pi(a)$ and $\pi(b)$, contradicting our assumption. 
\end{proof}

I think this proof is slightly more elegant:
\begin{lemma}\label{eqRelApproxFin}
  Suppose $\sim$ is a closed equivalence relation on a finite type $F$. Then $\sim$ can be written as a limit of decidable equivalence relations. 
\end{lemma}
\begin{remark}
  The above doesn't necessarily hold if we let $F$ be any profinite set. 
\end{remark}
\begin{proof}
  By \parencite[Thm 3.11.]{synthetic-stone-duality}, there is some $\alpha : \N \to F \to F \to 2$ be such that $a \sim b \leftrightarrow \forall_{n:\N} \alpha_n = 0$. 
  For $n:\N$ define $\sim_n$ as the equivalence relation induced by $R_n(a,b) := \forall_{k\leq n} \alpha_k(a,b) = 0$. As $F$ is finite, $\sim_n$ is decidable. It remains to be shown that $a\sim b \leftrightarrow \forall_{n:\N} a\sim_n b$. It is clear that if $a\sim b$ then $R_n(a,b)$ for all $n:\N$ and hence $a \sim_n b$ for all $n:\N$. 
  Conversely, by Markov's principle it is sufficient to show the contraposition. So suppose that $a \nsim_n b$ for some $n:\N$. Then in particular we cannot have $R_n(a,b)$ and thus there is some $k:\N$ with $\alpha_k(a,b) = 1$ and thus $a\nsim b$. 
\end{proof}

\begin{lemma}
  Let $F$ finite and $\sim $ a closed equivalence relation on $F$. 
  Then $X:=F/\sim$ is Stone. 
\end{lemma}
\begin{proof}
  By definition, $X$ is compact Hausdorff. It is then Stone if it is totally disconnected \parencite[Lemma 4.17]{synthetic-stone-duality}. We thus need to show 
  $$
  \Pi_{x,y:X} ((\Pi_{D:X\to 2} D(x) = D(y)) \to x = y)
  $$
  It is sufficient to show that if $a,b:F$ are such that $D(\pi(a)) = D(\pi(b))$ for all $D:F \to 2$ respecting $\sim$, then $a\sim b$. 
  As $a\sim b$ is closed, we can instead show the double negation, and by \Cref{eqRelApproxFin} assume that $\sim$ is the limit of decidable relations $\sim_n$. So suppose $a\nsim b$. By Markov there is some $k:\N$ with $a\nsim_k b$. Define $E:F\to 2$ by $E(c) = a\sim_k b$. $E$ respects $\sim$, and hence induces a map $X\to 2$ seperating $\pi(a)$ and $\pi(b)$, contradicting our assumption. 
\end{proof}




